\documentclass[10pt,
% a4paper,
%twocolumn,
fleqn,
%landscape, 
% papersize,
dvipdfmx,
uplatex
]{jsarticle}



\def\maru#1{\textcircled{\scriptsize#1}}%丸囲み番号

% \RequirePackage[2020/09/30]{platexrelease}

%太字設定
\usepackage[deluxe]{otf}

\usepackage{emathEy}

\usepackage[g]{esvect}

%大きな文字
\usepackage{fix-cm}

%定理環境
\usepackage{emathThm}
%\theoremstyle{boxed}
\theorembodyfont{\normalfont}
\newtheorem{Question}{問題}[subsection]
\newtheorem{Q}{}[subsection]
\newtheorem{question}[Question]{}
\newtheorem{quuestion}{}[subsection]

%セクション，大問番号のデザイン
\renewcommand{\labelenumi}{(\arabic{enumi})}
\renewcommand{\theenumii}{\alph{enumii})}
\renewcommand{\thesection}{第\arabic{section}章}

%用紙サイズの詳細設定
\usepackage{bxpapersize}
\papersizesetup{size={80mm,45mm}}
\usepackage[top=0.7zw,bottom=0truemm,left=3truemm,right=133truemm]{geometry}
\usepackage[dvipdfmx]{graphicx}

%余白など
\usepackage{setspace} % 行間
\setlength{\mathindent}{1zw}
\setlength\parindent{0pt}


%色カラーに関する設定
\usepackage{color}
\definecolor{shiro}{rgb}{0.95703125,0.87109375,0.7421875}
\definecolor{kin}{rgb}{0.95703125,0.87109375,0.7421875}
\definecolor{orange}{rgb}{1,0.7,0.2}
\definecolor{bradorange}{rgb}{1,0.5,0}
\definecolor{pink}{rgb}{0.9176,0.5686,0.5960}
\definecolor{mizu}{rgb}{0.6156,0.8,0.9955}
\color{kin}
% \pagecolor{hukamido}

\usepackage{at}%図の配置
% \usepackage{wallpaper}

\begin{document}



\at(0cm,0cm){\includegraphics[width=8cm,bb=0 0 1920 1080]{./youtube/thumbnails/templates/smart_background/図形と方程式.jpeg}}
{\color{orange}\bf\boldmath\Large\underline{円の方程式の基本形と一般形}}\vspace{0.3zw}

\large 
\bf\boldmath 問．次の方程式はどのような図形を表すか．\\
(1)  $x^2+y^2-2x+6y=0$\\
(2)  $x^2+y^2-6x+8y+{25}=0$\\
(3)  $x^2+y^2-2x+4y+{10}=0$\\

\at(6.6cm,0.2cm){\small\color{bradorange}$\overset{\text{図形と方程式}}{\text{基本}}$}


\newpage



\at(0cm,0cm){\includegraphics[width=8cm,bb=0 0 1920 1080]{./youtube/thumbnails/templates/smart_background/図形と方程式.jpeg}}
{\color{orange}\bf\boldmath\huge\underline{円の方程式}}\vspace{0.3zw}

\large 
\bf\boldmath 問．次の円の方程式を求めよ．\\
(1)  中心が$\left(-1,\;0\right)$で半径$2$の円\\
(2)  $2$点$\text{A}\left(3,\;4\right),\;\text{B}\left(-1,\;2\right)$を両端とする線分を直径とする円\\

\at(6.6cm,0.2cm){\small\color{bradorange}$\overset{\text{図形と方程式}}{\text{基本}}$}


\newpage



\at(0cm,0cm){\includegraphics[width=8cm,bb=0 0 1920 1080]{./youtube/thumbnails/templates/smart_background/図形と方程式.jpeg}}
{\color{orange}\bf\boldmath\huge\underline{線対称な点の座標}}\vspace{0.3zw}

\large 
\bf\boldmath 問．次のような点$\text{P}$の座標を求めよ．\\
(1)  点$\text{A}\left(2,\;1\right)$に関して，点$\text{B}\left(-2,\;3\right)$と対称な点$\text{P}$\\
(2)  直線$l：2x-y-1=0$に関して点$\text{B}\left(0,\;4\right)$と対称な点$\text{P}$\\

\at(6.6cm,0.2cm){\small\color{bradorange}$\overset{\text{図形と方程式}}{\text{基本}}$}


\newpage



\at(0cm,0cm){\includegraphics[width=8cm,bb=0 0 1920 1080]{./youtube/thumbnails/templates/smart_background/図形と方程式.jpeg}}
{\color{orange}\bf\boldmath\huge\underline{三角形の面積}}\vspace{0.3zw}

\LARGE 
\bf\boldmath 問．点$\text{A}\left(3,\;7\right),\;\text{B}\left(1,\;1\right),\;\text{C}\left(4,\;5\right)$を頂点とする三角形$\text{ABC}$の面積を求めよ．
\at(6.6cm,0.2cm){\small\color{bradorange}$\overset{\text{図形と方程式}}{\text{基本}}$}


\newpage



\at(0cm,0cm){\includegraphics[width=8cm,bb=0 0 1920 1080]{./youtube/thumbnails/templates/smart_background/図形と方程式.jpeg}}
{\color{orange}\bf\boldmath\LARGE\underline{点から直線に至る距離}}\vspace{0.3zw}

\Large 
\bf\boldmath 問．次の点と直線の距離を求めよ．\\
(1)  点$\left(2,\;-3\right)$と直線$2x+y-3=0$\\
(2)  点$\left(-1,\;5\right)$と直線$y=3x-2$\\

\at(6.6cm,0.2cm){\small\color{bradorange}$\overset{\text{図形と方程式}}{\text{基本}}$}


\newpage



\at(0cm,0cm){\includegraphics[width=8cm,bb=0 0 1920 1080]{./youtube/thumbnails/templates/smart_background/図形と方程式.jpeg}}
{\color{orange}\bf\boldmath\large\underline{$2$直線の傾きによる平行条件と垂直条件}}\vspace{0.3zw}

\LARGE 
\bf\boldmath 問．点$\left(3,\;1\right)$を通り，直線$3x+2y+1=0$に垂直な直線の方程式を求めよ．
\at(6.6cm,0.2cm){\small\color{bradorange}$\overset{\text{図形と方程式}}{\text{基本}}$}


\newpage



\at(0cm,0cm){\includegraphics[width=8cm,bb=0 0 1920 1080]{./youtube/thumbnails/templates/smart_background/図形と方程式.jpeg}}
{\color{orange}\bf\boldmath\huge\underline{直線の方程式}}\vspace{0.3zw}

\large 
\bf\boldmath 問．次の直線の方程式を求めよ．\\
(1)  点$\left(2,\;4\right)$を通り傾きが$3$\\
(2)  $2$点$\left(3,\;2\right),\;\left(5,\;6\right)$を通る\\
(3)  $2$点$\left(3,\;-1\right),\;\left(3,\;4\right)$を通る\\

\at(6.6cm,0.2cm){\small\color{bradorange}$\overset{\text{図形と方程式}}{\text{基本}}$}


\newpage



\at(0cm,0cm){\includegraphics[width=8cm,bb=0 0 1920 1080]{./youtube/thumbnails/templates/smart_background/図形と方程式.jpeg}}
{\color{orange}\bf\boldmath\huge\underline{$2$点間の距離}}\vspace{0.3zw}

\LARGE 
\bf\boldmath 問．点$\text{P}\left(2,\;1\right)$からの距離が$\sqrt {10}$である$x$軸上の点$\text{Q}$の座標を求めよ．
\at(6.6cm,0.2cm){\small\color{bradorange}$\overset{\text{図形と方程式}}{\text{基本}}$}


\newpage



\at(0cm,0cm){\includegraphics[width=8cm,bb=0 0 1920 1080]{./youtube/thumbnails/templates/smart_background/図形と方程式.jpeg}}
{\color{orange}\bf\boldmath\huge\underline{三角形の重心の座標}}\vspace{0.3zw}

\normalsize
\bf\boldmath 点$\text{A}\left(a_1,\;a_2\right),\;\text{B}\left(b_1,\;b_2\right),\;\text{C}\left(c_1,\;c_2\right)$を頂点とする三角形$\text{ABC}$の重心$\text{G}$の座標は，

\Large
\vspace{0.4zw}
\hspace{0zw}
$\left(\bunsuu{a_1+b_1+c_1}{3},\bunsuu{a_2+b_2+c_2}{3}\right)$

\large
\vspace{0.4zw}
\hfill
である．
\at(6.6cm,0.2cm){\small\color{bradorange}$\overset{\text{図形と方程式}}{\text{基本}}$}



\newpage



\at(0cm,0cm){\includegraphics[width=8cm,bb=0 0 1920 1080]{./youtube/thumbnails/templates/smart_background/図形と方程式.jpeg}}
{\color{orange}\bf\boldmath\huge\underline{外分点の座標}}\vspace{0.3zw}

\large
2点A$(a)$，B$(b)$を$m:n$に外分する点Pの座標は

\Huge
\hspace{1zw}
$\bunsuu{-na+mb}{m-n}$

\large
\vspace{-1zw}
\hfill
である．

\at(6.6cm,0.2cm){\small\color{bradorange}$\overset{\text{図形と方程式}}{\text{基本}}$}


\newpage



\at(0cm,0cm){\includegraphics[width=8cm,bb=0 0 1920 1080]{./youtube/thumbnails/templates/smart_background/図形と方程式.jpeg}}
{\color{orange}\bf\boldmath\huge\underline{内分点の座標}}\vspace{0.3zw}

2点A$(a)$，B$(b)$を$m:n$に内分する点Pの座標は

\Huge
\hspace{1.4zw}
$\bunsuu{na+mb}{m+n}$

\large
\vspace{-1zw}
\hfill
である．

\at(6.6cm,0.2cm){\small\color{bradorange}$\overset{\text{図形と方程式}}{\text{基本}}$}


\newpage



\at(0cm,0cm){\includegraphics[width=8cm,bb=0 0 1920 1080]{./youtube/thumbnails/templates/smart_background/図形と方程式.jpeg}}
{\color{orange}\bf\boldmath\huge\underline{円の接線の公式}}\vspace{0.3zw}

\vspace{0.4zw}
円$x^2+y^2=r^2$の点$(p,q)$における接線の方程式は

\HUGE
\vspace{-0.4zw}
\hspace{0.2zw}
$px+qy=r^2$

\large
\vspace{0.5zw}
\hfill
である．
\at(6.6cm,0.2cm){\small\color{bradorange}$\overset{\text{図形と方程式}}{\text{基本}}$}


\end{document}

